% !TEX root = ../main.tex
\section{Sub-Doxastic Systems and the Limits of the Analogy}
\label{sec:subdoxastic}

\subsection{The positive characterisation}

If machine assertion-states are not beliefs, what are they?
\citet{stich1978} introduced the category of the \emph{sub-doxastic} for states
that carry information and influence behaviour while failing two marks of belief:
they are not inferentially integrated with the agent's other states, interacting
with beliefs and with each other only along narrow special-purpose channels; and
they are not available at the person level. His examples are the states of the
visual system computing depth from disparity, and the internalised grammar
generating acceptability intuitions. A speaker's parser registers that a string
is ill-formed, but this registration does not enter general reasoning the way a
belief would: it cannot be weighed against other evidence, deployed in planning,
or revised by argument.

The category is the right one for these systems, but it needs reconstruction on
both of Stich's marks rather than one, and I would rather do that openly than
borrow the term and hope the differences go unnoticed.

\subsection{Why inferential isolation is the wrong mark}
\label{sec:isolation}

It is tempting to say that machine assertion-states satisfy Stich's
isolation criterion superbly, and to read the axiom failures of
Section~\ref{sec:output-operator} as its formal signature. That reading conflates
two properties, and on the one Stich actually specified, these systems fail the
criterion---in an unexpected direction.

Stich's isolation is \emph{encapsulation}: a boundary condition. The parser's
states cannot interact with the general belief set. What the normal axioms
formalise is something else: \emph{closure}, a discipline condition on a set of
contents. These are independent. A module may be internally well-disciplined and
informationally sealed, which is roughly the picture of the modular mind that
Stich's examples belong to. And a system may be radically unencapsulated and
wholly undisciplined: everything interacting with everything, just not logically.

The second description fits a language model better than the first. Content in a
forward pass is globally available to the computation; there is no wall, and
nothing is special-purpose. What these systems exhibit is not inferential
isolation but inferential \emph{promiscuity without discipline}---which is, if
anything, the opposite of encapsulation. So the analogy does not survive on this
mark, and the appearance that it does comes from hearing ``isolation'' as
``failure to draw the right conclusions'' rather than as ``walled off.''

\subsection{Why the personal-level relatum is missing}
\label{sec:relatum}

The second mark fails for a different and more interesting reason. Stich's notion
is defined relationally. Sub-doxastic states are \emph{sub-personal}: states of
subsystems within a creature that also has a personal level, a locus of belief,
inference, and rational agency relative to which they are ``sub.'' The visual
system's disparity-states are sub-doxastic of a perceiver; the parser's states are
sub-doxastic of a speaker. The inferential isolation that defines them is
isolation \emph{from the agent's beliefs}.

In the machine case there is no personal-level believer for these states to be
sub-personal relative to. The system is, so to speak, sub-doxastic all the way
up: a parser without a speaker, a visual module without a perceiver. This is not
a quibble. It removes the contrast that gave Stich's category its content.
Strictly, then, these states are not sub-doxastic in Stich's sense, because
nothing satisfies the relatum his definition requires.

\subsection{The monadic core}
\label{sec:monadic}

What survives both departures is a monadic core, which the apparatus of this
paper lets us state without the missing relatum and without the misdescribed
mark. Call a content-bearing state sub-doxastic in the extended sense iff it is:

\begin{enumerate}
  \item \textbf{informational}---apt for grounding in an evidential base, hence
        apt for $\Grd_M$;
  \item \textbf{behaviour-guiding}---issuing in assertion, hence apt for
        $\Out_M$; and
  \item \textbf{closure-free}---failing the normal axioms constitutively rather
        than by performance limitation.
\end{enumerate}

Call systems all of whose content-bearing states meet (1)--(3) \emph{sub-doxastic
systems} in an extended sense: not systems containing sub-doxastic states beneath
a doxastic level, but systems exhibiting the sub-doxastic profile with no doxastic
level anywhere. The extension is a genuine conceptual novelty and I flag it as
such. Human cognitive science never needed the monadic notion, because in humans
sub-doxastic states always come embedded in a believer. These may be the first
well-studied free-standing sub-doxastic systems.

Condition (3) is where the philosophical weight sits, and it is worth connecting
to why closure is a norm for belief in the first place. Belief is constitutively
aimed at truth \citep{shah2005,wedgwood2002}; assertion is governed by a norm
that holds the asserter to what she is in a position to know
\citep{williamson2000}, and to the commitments her assertion incurs
\citep{brandom1994}. Closure is not an ornament on belief but a consequence of
its aim: a state that will be recruited by open-endedly many projects must be
kept consistent and closed, on pain of the projects failing. A state that carries
information and guides output while answering to no such aim is exactly what
conditions (1)--(3) describe.

The characterisation earns its keep by steering between the two crude positions
of Section~\ref{sec:intro}. Against the enthusiast: sub-doxastic states are not
beliefs, and no scaling of the same architecture makes them beliefs, because the
deficit is structural. Against the eliminativist: sub-doxastic states are not
nothing---they carry information, admit of grounding and ungrounding, and support
a genuine if non-doxastic epistemology of the kind
Sections~\ref{sec:grounding}--\ref{sec:hallucination} provide. ``Stochastic
parrot'' is as much an underdescription as ``artificial believer.''

Two corollaries. First, familiar person-level appraisal---rationality,
irrationality, dishonesty, sincerity---misfires for such systems not because they
are bad believers but because there is no believer; the appraisals lack a
subject. Second, the appraisals that do apply are the mechanism-level ones of
Section~\ref{sec:hallucination}: grounded or ungrounded, faithful or fabricating,
sensitive or lucky. The sub-doxastic framing and the hallucination framework are
again two faces of one view.
